\documentclass[10pt,a4paper]{amsart}
\usepackage[margin=19mm]{geometry}
\usepackage{amsmath,amssymb,mathtools}
\usepackage[scr=rsfs,cal=euler]{mathalfa}
\usepackage{xfrac}
\usepackage[unicode,hidelinks,bookmarks=false]{hyperref}
\newcommand{\ie}{i.e.\ }
\newcommand{\cf}{cf.\ }
\newcommand{\resp}{respectively\ }
\newcommand{\Subsection}{\subsection}
\providecommand{\nobreakdash}{\nobreak\hbox{-}}
\setlength{\emergencystretch}{2em}
\multlinegap0pt
\newtheorem{theorem}{Theorem}
\newtheorem{proposition}{Proposition}
\DeclareMathOperator{\cst}{constant}
\DeclareMathOperator{\trace}{trace} \DeclareMathOperator{\di}{div}
\title{Simons formulas in complex space forms and product spaces}
\author{Dorel Fetcu, Stefano Nardulli}
\date{Source: arXiv:2608.14238v1 (CC BY 4.0)}
\begin{document}
\maketitle

\noindent\emph{Source and licence.} Simons formulas in complex space forms and product spaces. Version arXiv:2608.14238v1 (CC BY 4.0), distributed by arXiv under the Creative Commons Attribution 4.0 International licence, http://creativecommons.org/licenses/by/4.0/, as recorded in the arXiv licence metadata for this exact version and shown on the abstract page https://arxiv.org/abs/2608.14238v1. Full licence notice: \texttt{licenses/arxiv-2608.14238-CC-BY-4.0.txt}.\par\smallskip

\noindent\textit{Editorial scope.} Counted unit: the article's proposition hyper (source0.tex lines 810--815). The article's complete original proof of it is delivered with it (source0.tex lines 817--843, one \texttt{proof} environment of 27 lines). No mathematics is written, repaired or re-derived: the statement and the proof are the article's own lines, in the article's own notation, and no retained line is substituted or re-flowed.  Retained uncounted as the source-local context the proof needs: lines 793--799.  Nothing was omitted from any retained span, so the package carries no omission record.  No retained line contains a citation, so the wrapper carries no bibliography and no bibliography entry.  No retained line contains a figure environment, an image-inclusion command or drawing code, so no image is delivered and the package declares no asset dependency.  Editorial formatting only, in addition: an AMS \texttt{amsart} preamble that restates the article's own declarations and macros used by the retained lines, namely the environments \texttt{theorem}, \texttt{proposition} on separate counters, together with the standard packages those lines need.  The retained environments print in this document as Theorem 1, Proposition 1 in order of appearance, whereas the article prints the same items with the numbering of its own full document.  The complete native source of the article as distributed in the arXiv e-print for this exact version is bundled unchanged as \texttt{sources/207648/source0.tex}.
\begin{theorem}\label{thm:delta2_prod}
Let $\Sigma^m$ be a submanifold of $N^n(c)\times\mathbb{R}$. If $V$ is a normal vector field parallel in the normal bundle, with $\trace A_V=\cst$, then
\begin{eqnarray}\label{eq:delta2_prod}
\frac{1}{2}\Delta|A_V|^2&=&|\nabla A_V|^2+\frac{c}{4}\{m|A_V|^2-2m|A_VT|^2+6|tA_V|^2-|T|^2|A_V|^2\\\nonumber &&+3(\trace A_V)\langle A_VT,T\rangle+6\trace(A_V(tA_{nV}))\\\nonumber &&+3\trace(A_VA_{n(tV)})-(\trace A_V)^2+6\trace(A_V(t(A_Vt)))\\\nonumber &&+3\trace(\langle t(\sigma(A_V\cdot,tV)),\cdot\rangle)-3\trace(\langle t(\sigma(A_V(tV),\cdot),\cdot\rangle)\\\nonumber &&+m\langle V,N\rangle(\trace(A_VA_N)-(\trace A_V)\langle H,N\rangle)\\\nonumber&&+3m\langle A_V(tH),tV\rangle\}+\sum_{\alpha=2m+1}^{2n}\{\trace(A_VA_{\alpha}A_VA_{\alpha}-A_V^2A_{\alpha}^2)\\\nonumber &&+\trace(A_{\alpha})\trace(A_V^2A_{\alpha})-(\trace(A_VA_{\alpha}))^2\},
\end{eqnarray}
where $\{E_{\alpha}\}_{\alpha=m+1}^{n+1}$ is a local orthonormal frame field in the normal bundle.
\end{theorem}

\begin{proposition}\label{hyper} If $\Sigma^{2n}$ is a CMC hypersurface in $N^n(c)\times\mathbb{R}$, then
\begin{eqnarray}\label{eq:hyper2_prod}
\frac{1}{2}\Delta|A|^2&=&|\nabla A|^2-|A|^4+2n|H|\trace A^3+\frac{c}{4}\{(4n-(2n+4)|T|^2)|A|^2\\\nonumber &&-4n|AT|^2+3|tA-At|^2+6n|H|\langle AT,T\rangle\\\nonumber&&+6n|H|\langle A(\varphi\nu),\varphi\nu\rangle-4n^2(2-|T|^2)|H|^2\}
\end{eqnarray}

\end{proposition}

\begin{proof} If at a point $q\in\Sigma$ we have $T_q=0$, then one can see that \eqref{eq:hyper2_prod} actually turns out to be \eqref{eq:delta2_prod}, at $q$. Let us assume that $T\neq 0$  at a point $p$ on the surface. Then, at $p$, the vector fields $\varphi\nu$ and $T$ are tangent and orthogonal to each other. Moreover, $|\varphi\nu|^2=1-\eta^2(\nu)=|T|^2$. Thus we can define $E_1=-(\varphi\nu)/|T|$ and $E_2=T/|T|$, two unit tangent vectors. Now, if $E_3$ is another unit tangent vector orthogonal to $E_1$ and $E_2$, it follows that
$$
\langle\varphi E_3,\nu\rangle=-\langle E_3,\varphi\nu\rangle=0\quad\textnormal{and}\quad \langle\varphi E_3,\varphi E_3\rangle=\langle E_3,E_3\rangle,
$$
since $E_3\perp\xi$, which means that $\varphi E_3$ is also a unit tangent vector, and then
$$
\langle\varphi E_3,E_1\rangle=-\frac{1}{|T|}(\langle E_3,\nu\rangle-|N|\langle E_3,\xi\rangle)=0,\quad\langle\varphi E_3,E_2\rangle=\frac{1}{|T|}\langle\varphi E_3,\xi\rangle=0.
$$ 
Therefore we can consider an orthonormal basis
$$
\{E_1,E_2,E_3,E_4=\varphi E_3\ldots,E_{2n-1},E_{2n}=\varphi E_{2n-1}\},
$$
in $T_p\Sigma$. It is also easy to compute 
$$
\varphi E_1=-|N|E_2+|T|\nu\quad\textnormal{and}\quad\varphi E_2=|N|E_1.
$$

Next, for any $i\geq 2$, we have $\langle t(\sigma(AE_i,\varphi\nu),E_i\rangle=\langle t(\sigma(E_i,A(\varphi\nu)),E_i\rangle=0$. Moreover, we also obtain 
$$
\langle t(\sigma(AE_1,\varphi\nu),E_1\rangle=\langle t(\sigma(E_1,A(\varphi\nu)),E_1\rangle=\langle\sigma(AE_1,E_1),\nu\rangle=|A(\varphi\nu)|^2
$$
and $|tA|^2=|At|^2=|A|^2-|AE_1|^2-|AE_2|^2$, which implies that 
$$
|tA|^2+\trace(A(t(At)))=(1/2)|tA-At|^2.
$$
Since $n\nu=0$ and $n(t\nu)=n(\varphi\nu)=-|T|^2\nu$, we replace in equation \eqref{eq:delta2_prod} and conclude the proof.
\end{proof} 

\end{document}
